\Askhsh{13672}{4}
\begin{ylh}{1}
    \item 3.1. Είδη και στοιχεία τριγώνων
    \item 3.2. 1ο Κριτήριο ισότητας τριγώνων
    \item 4.2. Τέμνουσα δύο ευθειών - Ευκλείδειο αίτημα
    \item 5.9. Μια ιδιότητα του ορθογώνιου τριγώνου.
\end{ylh}
% ===================================================================
%  Αρχείο: 13672.tex (Fragment)
%  Αυτόματη μετατροπή μέσω doc2latex.py
% ===================================================================

ΘΕΜΑ 4

Δίνεται ορθογώνιο τρίγωνο ΑΒΓ με $\widehat{\text{Α}} = 90^{ο}$ και $\text{ΑΒ} >$ $\text{ΑΓ}$. Από το μέσο Δ της πλευράς ΒΓ φέρουμε κάθετη στη ΒΓ όπως φαίνεται στο παρακάτω σχήμα, η οποία τέμνει τη διχοτόμο ΑΗ της γωνίας $\widehat{\text{Α}}$ στο σημείο Ε. Έστω ΑΖ το ύψος στην υποτείνουσα. Να αποδείξετε ότι:

\begin{alist}
\item $\text{Γ}\widehat{\text{Α}}\text{Ζ} =$ $\text{Δ}\widehat{\text{Α}}\text{Β}$.

\monades{8}

\item $\text{ΑΔ} = \text{ΔΕ}$.

\monades{9}

\item $\text{Ζ}\widehat{\text{Α}}\text{Δ} = \widehat{\text{Γ}} - \widehat{\text{Β}}$.

\monades{8}

\includesvg[width=2.39055in,height=2.71654in]{/home/spyros/Μαθηματικά/Trapeza_Thematwn/A_Lykeiou/Geometry/extracted/13672/media/image2.svg}
\end{alist}
